% 제출본 전용 — 10장 끝에 붙는다 (10.5). 공학인증 현실적 제약 조건 표.

\subsection{현실적 제약 조건}

설계는 제약 안에서 이루어진다. 이 작품이 어떤 제약을 어떻게 다뤘는지를 정리한다.
10.4절의 한계가 \textbf{결과로서 말할 수 없는 것}이라면, 이 절은 \textbf{설계
과정에서 감당해야 했던 조건}이다.

\begin{table}[htbp]
  \centering
  \caption{현실적 제약 조건}
  \label{tab:v2-constraint}
  \begin{tabular}{llcp{8.6cm}}
    \toprule
    \textbf{구분} & & \textbf{연관성} & \textbf{내용} \\
    \midrule
    \multirow{6}{*}{\parbox{1.6cm}{현실적\\제한요소}}
      & 비용 & O &
        \textbf{추가 비용 0원.} 공개 자료(국토교통부 실거래 API, 서울시보·
        전자관보 공고)와 개인 PC 만 쓴다. 유료 시세 자료를 쓰지 않아도
        답할 수 있는 범위로 질문을 정했다 \\
    \cmidrule(l){2-4}
      & 기간 & O &
        \textbf{1인 12주.} M1 주 분석 확정 --- M2 하이라이트 --- M3 분석
        종료 --- M4 최종본으로 쪼개고 \textbf{매주 보고서}로 점검했다
        (1.5절). 1인이므로 병렬 작업이 불가능해, 확장 과제는 처음부터
        범위 밖으로 두었다 \\
    \cmidrule(l){2-4}
      & 정확성 & O &
        \textbf{현재 값의 출처를 한 곳으로 묶었다}(0007). 보고서의 표와
        결과 대시보드는 \texttt{output/results/*.json}에서 생성되고,
        문서의 숫자가 그 파일과 어긋나면 매 실행에서 걸린다. 실제로 석 달
        묵은 값과 여섯 달 묵은 값을 이 장치가 잡았다(3.7절) \\
    \cmidrule(l){2-4}
      & 신뢰성 & O &
        처치 단위가 적어 관행적 $p$ 값을 쓸 수 없으므로 \textbf{무작위화
        추론과 와일드 부트스트랩}으로 바꿨다(5.3절). 평행추세 가정도
        그대로 믿지 않고 \textbf{풀어 가며} 결론이 무너지는 지점을 쟀다
        (8.4절). \textbf{가릴 수 없는 것은 가릴 수 없다고 보고한다} \\
    \cmidrule(l){2-4}
      & 안전성 & O &
        API 인증키를 저장소에 넣지 않고 \texttt{.env}(git 제외)에서 읽는다.
        분석의 최소 단위는 (단지, 전용면적, 층)이며 \textbf{특정 호실이
        아니므로 개인을 식별하지 않는다}. 원자료 902\,MB 와 서울시보 원본
        2.9\,GB 도 저장소 밖에 둔다 (부록 B) \\
    \cmidrule(l){2-4}
      & \parbox{2.2cm}{사회에\\미치는 영향} & O &
        정책 효과 추정은 해석이 오남용되기 쉽다. 그래서 \textbf{모든 효과를
        반사실 대비로} 쓰고(0006), 숫자마다 조건을 붙이며(0007), 일반화
        범위를 \textbf{``서울 고가 재건축 지역''까지}로 명시한다(10.4절).
        가격 효과는 크기를 쓰지 않고 \textbf{방향만} 쓴다(8.6절) \\
    \bottomrule
  \end{tabular}
\end{table}

\textbf{마지막 줄이 이 작품에서 가장 조심한 부분이다.} 토지거래허가제는 지금도
시행 중인 제도이고, ``규제가 가격을 $x$\% 낮췄다''는 문장은 인용되기 쉽다. 그런데
8장에서 보였듯 가격 효과는 평행추세 가정을 조금만 풀면 0과 구분되지 않는다.
\textbf{말할 수 있는 것과 없는 것을 가르는 일이 이 작품에서는 결과를 내는 일만큼
중요했다.}
